% ---- pg-456 （书页 444） ----

\noindent 10.13\quad 沃(Waugh, 见文献[12])提供了容量为 $n=138$ 的加拿大硬红春麦样本的信息, 以及用该样本磨成的面粉的信息. $p=5$ 个小麦测度量(标准化形式)为

$\mathbf{z}_1^{(1)}=$ 谷粒构造

$\mathbf{z}_2^{(1)}=$ 试验重量

$\mathbf{z}_3^{(1)}=$ 破损的谷粒

$\mathbf{z}_4^{(1)}=$ 异地原料

$\mathbf{z}_5^{(1)}=$ 小麦中的天然蛋白质

$q=4$ 个(标准化)面粉测度量为

$\mathbf{z}_1^{(2)}=$ 磨成每桶面粉需要小麦量

$\mathbf{z}_2^{(2)}=$ 面粉中的灰

$\mathbf{z}_3^{(2)}=$ 面粉中的天然蛋白质

$\mathbf{z}_4^{(2)}=$ 麦麸的品质指标

样本相关系数矩阵为
\[
\mathbf{R}=\left[\begin{array}{c:c}\mathbf{R}_{11}&\mathbf{R}_{12}\\ \hdashline \mathbf{R}_{21}&\mathbf{R}_{22}\end{array}\right]=\left[\begin{array}{ccccc:cccc}
1.0& & & & & & & & \\
0.754&1.0& & & & & & & \\
-0.690&-0.712&1.0& & & & & & \\
-0.446&-0.515&0.323&1.0& & & & & \\
0.692&0.412&-0.444&-0.334&1.0& & & & \\ \hdashline
-0.605&-0.722&0.737&0.527&-0.383&1.0& & & \\
-0.479&-0.419&0.361&0.461&-0.505&0.251&1.0& & \\
0.780&0.542&-0.546&-0.393&0.737&-0.490&-0.434&1.0& \\
-0.152&-0.102&0.172&-0.019&-0.148&0.250&-0.079&-0.163&1.0
\end{array}\right]
\]

\noindent (a)\quad 求与显著的(在 $\alpha=0.01$ 水平)典型相关系数相对应的样本典型变量.

\noindent (b)\quad 解释第一样本典型变量 $\hat{U}_1,\hat{V}_1$. 从某种意义上说, 它们能否分别代表整体的小麦和面粉质量?

\noindent (c)\quad 第一组 $\mathbf{Z}^{(1)}$ 的样本总方差被典型变量 $\hat{U}_1$ 解释的比例为多少? 第二组 $\mathbf{Z}^{(2)}$ 的样本总方差被典型变量 $\hat{V}_1$ 解释的比例为多少? 讨论一下你的答案.

\noindent 10.14\quad 考虑练习 9.15 给出的盈利测度的相关系数矩阵. 令 $\mathbf{X}^{(1)}=[X_1^{(1)},X_2^{(1)},\cdots,X_p^{(1)}]'$ 代表盈利的会计测度向量, 而令 $\mathbf{X}^{(2)}=[X_1^{(2)},X_2^{(2)}]'$ 为盈利的两个市场测度向量.

将样本相关系数矩阵分块并作典型相关分析, 特别地:

\noindent (a)\quad 确定第一样本典型变量 $\hat{U}_1,\hat{V}_1$ 及它们的相关系数. 解释这些典型变量.

\noindent (b)\quad 令 $\mathbf{Z}^{(1)}$ 和 $\mathbf{Z}^{(2)}$ 分别为对应于 $\mathbf{X}^{(1)}$ 和 $\mathbf{X}^{(2)}$ 的标准化变量组. $\mathbf{Z}^{(1)}$ 的样本总方差用典型变量 $\hat{U}_1$ 解释的比例为多少? $\mathbf{Z}^{(2)}$ 的样本总方差用典型变量 $\hat{V}_1$ 解释的比例又是多大? 请讨论.

\noindent 10.15\quad 四项关于刚性的测度值, 在\hyperref[tab:4.3]{表 4.3} 中给出并在例 4.14 中讨论过, 用表中数据构造样本协方差矩阵 $\mathbf{S}$. 令向量 $\mathbf{X}^{(1)}=[X_1^{(1)},X_2^{(1)}]'$ 的变量代表刚性动态测度(冲击波、波动), 并
